%%% FIELD proof-prop-two-arm-reduction
Write the distinct atoms of \(F_0^{(2)}\) and \(F_1^{(2)}\) as \(\alpha_1,\ldots,\alpha_R\) and \(\beta_1,\ldots,\beta_S\), respectively, and put
\[
k_r=NF_0^{(2)}(\{\alpha_r\}),\qquad
\ell_s=NF_1^{(2)}(\{\beta_s\}).
\]
Because both laws are empirical laws of \(N\) fixed entries, every \(k_r\) and \(\ell_s\) is a positive integer and
\(\sum_r k_r=\sum_s\ell_s=N\).  Every coupling in \(\Gamma(F^{(2)})\) is supported on the finite Cartesian product of these atom sets.  Thus it is represented by a nonnegative matrix \(p=(p_{rs})\) satisfying
\[
\sum_s p_{rs}=\frac{k_r}{N},\qquad
\sum_r p_{rs}=\frac{\ell_s}{N}.
\]
After the invertible scaling \(q=Np\), the feasible set is the transportation polytope
\[
q_{rs}\geq 0,\qquad \sum_s q_{rs}=k_r,\qquad \sum_r q_{rs}=\ell_s. \tag{1}
\]

We first prove the required integrality rather than assume a transportation-polytope theorem.  Associate with a feasible \(q\) the bipartite graph having a row vertex \(r\), a column vertex \(s\), and an edge \((r,s)\) exactly when \(q_{rs}>0\).  If this positive-support graph contains a cycle, assign alternating signs \(+1,-1,+1,-1,\ldots\) to its edges and zero off the cycle.  The resulting matrix \(h\neq0\) has every row and column sum equal to zero.  Taking
\(0<\varepsilon<\min\{q_{rs}:(r,s)\text{ lies on the cycle}\}\) gives two distinct feasible matrices \(q+\varepsilon h\) and \(q-\varepsilon h\) whose midpoint is \(q\).  Hence the support of every extreme point is a forest.

Now let \(q\) be an extreme point.  In each nonempty component of its support forest, choose a leaf.  The weight on the unique edge incident to that leaf equals the leaf's residual margin in (1), which is an integer.  Delete that edge and leaf and subtract this integer weight from the neighboring residual margin.  The new residual margins remain nonnegative integers.  Iterating, while discarding any isolated zero-margin vertex, shows that every positive entry of \(q\) is an integer.  Expanding the cell \((r,s)\) into \(q_{rs}\) rows \((\alpha_r,\beta_s)\) produces exactly \(N\) rows and uses \(\alpha_r\) exactly \(k_r\) times and \(\beta_s\) exactly \(\ell_s\) times.  It therefore produces a member of \(\mathcal S_N(F^{(2)})\).  This proves that every extreme point of \(\Gamma(F^{(2)})\) is realizable.

The polytope (1) is nonempty, because \(q_{rs}=k_r\ell_s/N\) is feasible, and it is closed and bounded in a finite-dimensional space.  The linear cost therefore has a minimizer and a maximizer.  Starting from either optimizer, if its positive support contains a cycle, use the alternating perturbation above and increase \(\varepsilon\) until an edge first becomes zero.  Linearity says that the original objective is the average of the two perturbed objectives.  At a minimum both perturbed objectives are no smaller, so both equal the original objective; at a maximum the reversed inequalities give the same conclusion.  Retaining a perturbed optimizer with one fewer positive cell and iterating produces an acyclic, hence integral, optimizer.  Thus both extrema may be computed over \(\mathcal S_N(F^{(2)})\).  It remains to identify them.  Since \(d_0+d_1=0\), write \(d=(\lambda,-\lambda)\), where the case \(\lambda=0\) is included.  Expand the two empirical marginals into ordered lists
\[
A_1\leq\cdots\leq A_N,\qquad B_1\leq\cdots\leq B_N.
\]
After using the first list as row labels, every realizable table has cost
\[
\frac1N\sum_{i=1}^N(\lambda A_i-\lambda B_{\sigma(i)})^2
=\frac{\lambda^2}{N}\left(\sum_i A_i^2+\sum_iB_i^2-2\sum_iA_iB_{\sigma(i)}\right) \tag{2}
\]
for a permutation \(\sigma\).  If \(i<j\), \(A_i\leq A_j\), and the two paired values satisfy \(B_{\sigma(i)}>B_{\sigma(j)}\), swapping those two values changes the product sum by
\[
(A_iB_{\sigma(j)}+A_jB_{\sigma(i)})
 -(A_iB_{\sigma(i)}+A_jB_{\sigma(j)})
=(A_j-A_i)(B_{\sigma(i)}-B_{\sigma(j)})\geq0. \tag{3}
\]
Successively removing inversions proves that \(\sum_iA_iB_{\sigma(i)}\) is maximized by like ordering.  Applying the same exchange argument to \(-B_i\), or equivalently successively removing like-ordered pairs, proves that it is minimized by reverse ordering.  Equalities in (3) show that repeated atoms cause no problem.  Since each empirical quantile is constant on intervals of length \(1/N\), equation (2) now gives
\[
\min_{\pi\in\Gamma(F^{(2)})}\int(d_0y_0+d_1y_1)^2\,d\pi
=\int_0^1[d_0Q_0^{(2)}(u)+d_1Q_1^{(2)}(u)]^2\,du
=C_{\mathrm{co}}^{(2)}(F^{(2)}), \tag{4}
\]
and, with identities at the finitely many quantile boundaries irrelevant to the integral,
\[
\max_{\pi\in\Gamma(F^{(2)})}\int(d_0y_0+d_1y_1)^2\,d\pi
=\int_0^1[d_0Q_0^{(2)}(u)+d_1Q_1^{(2)}(1-u)]^2\,du
=C_{\mathrm{ctr}}^{(2)}(F^{(2)}). \tag{5}
\]
The same extrema over \(\mathcal S_N(F^{(2)})\) follow from integrality.

We next derive the design variance.  Fix any table \(T\in\mathcal S_N(F^{(2)})\), and let
\[
I_{bi}=\mathbf 1\{Z_i^{(2)}=b\},\qquad
S_{10}^{(2)}(T)=\frac1{N-1}\sum_{i=1}^N(T_{1i}-\mu_1^{(2)})(T_{0i}-\mu_0^{(2)}).
\]
The complete-randomization assumption and positivity of \(m_0,m_1\) give
\[
\mathbb E I_{bi}=\frac{m_b}{N},\qquad
\mathbb E(I_{bi}I_{bj})=\frac{m_b(m_b-1)}{N(N-1)}\quad(i\ne j), \tag{6}
\]
and
\[
\mathbb E(I_{0i}I_{1i})=0,\qquad
\mathbb E(I_{0i}I_{1j})=\frac{m_0m_1}{N(N-1)}\quad(i\ne j). \tag{7}
\]
In particular,
\[
\mathbb E_{Z^{(2)}}\widehat\theta_N^{(2)}
=\sum_{b\in\mathbb B}\frac{d_b}{m_b}\sum_i\frac{m_b}{N}T_{bi}
=\sum_{b\in\mathbb B}d_b\mu_b^{(2)}
=\theta_N^{(2)}, \tag{8}
\]
so the estimator is design-unbiased for the fixed causal contrast.

For completeness, center \(T_{bi}\) at \(\mu_b^{(2)}\).  Using (6), the facts that the centered entries sum to zero and that their squared sum is \((N-1)S_b^{2,(2)}\) yield
\[
\operatorname{Var}\!\left(\frac1{m_b}\sum_i I_{bi}T_{bi}\right)
=\left(\frac1{m_b}-\frac1N\right)S_b^{2,(2)}. \tag{9}
\]
Likewise, (7) and
\[
\sum_{i\ne j}(T_{0i}-\mu_0^{(2)})(T_{1j}-\mu_1^{(2)})
=-\sum_i(T_{0i}-\mu_0^{(2)})(T_{1i}-\mu_1^{(2)})
\]
give
\[
\operatorname{Cov}\!\left(\frac1{m_0}\sum_iI_{0i}T_{0i},
\frac1{m_1}\sum_iI_{1i}T_{1i}\right)
=-\frac1N S_{10}^{(2)}(T). \tag{10}
\]
Combining (9)--(10),
\[
\operatorname{Var}_{Z^{(2)}}(\widehat\theta_N^{(2)})
=\sum_{b\in\mathbb B}\frac{d_b^2S_b^{2,(2)}}{m_b}
-\frac1N\left\{\sum_{b\in\mathbb B}d_b^2S_b^{2,(2)}
+2d_0d_1S_{10}^{(2)}(T)\right\}. \tag{11}
\]
Define the fixed row contrast \(\tau_i=d_0T_{0i}+d_1T_{1i}\) and its average squared value
\(C_T=N^{-1}\sum_i\tau_i^2\).  Its mean is \(\theta_N^{(2)}\), because the marginal means do not depend on \(T\).  Therefore
\[
\frac1{N-1}\sum_i(\tau_i-\theta_N^{(2)})^2
=\frac{N}{N-1}\{C_T-(\theta_N^{(2)})^2\}
=\sum_b d_b^2S_b^{2,(2)}+2d_0d_1S_{10}^{(2)}(T). \tag{12}
\]
Since \(S_b^{2,(2)}=Nv_b^{(2)}/(N-1)\), substituting (12) into (11) gives the exact cost representation
\[
\operatorname{Var}_{Z^{(2)}}(\widehat\theta_N^{(2)})
=\frac{1}{N-1}\left\{
\sum_{b\in\mathbb B}\frac{N}{m_b}d_b^2v_b^{(2)}
+(\theta_N^{(2)})^2-C_T\right\}
=\frac{B_N^{(2)}(F^{(2)})-C_T}{N-1}. \tag{13}
\]

The first term in the numerator of (13) depends only on the fixed marginals.  Thus (4)--(5) imply that the maximum design variance over \(\mathcal S_N(F^{(2)})\) is obtained by the comonotone table and the minimum by the countermonotone table.  The schedule class is finite, so both extrema are attained; hence the interval hull is exactly
\[
\left[
\frac{B_N^{(2)}(F^{(2)})-C_{\mathrm{ctr}}^{(2)}(F^{(2)})}{N-1},
\frac{B_N^{(2)}(F^{(2)})-C_{\mathrm{co}}^{(2)}(F^{(2)})}{N-1}
\right]
=\bigl[L_N^{(2)}(F^{(2)}),U_N^{(2)}(F^{(2)})\bigr], \tag{14}
\]
where the last equality is exactly \(\mathrm{def:two-arm-endpoints}\).

Finally, expanding the comonotone cost around the two marginal means gives
\[
\frac{N}{N-1}\left\{C_{\mathrm{co}}^{(2)}(F^{(2)})-(\theta_N^{(2)})^2\right\}
=\sum_b d_b^2S_b^{2,(2)}+2d_0d_1\overline S_{10}^{(2)}, \tag{15}
\]
because
\(\overline S_{10}^{(2)}=\frac{N}{N-1}\{\int_0^1Q_1^{(2)}(u)Q_0^{(2)}(u)\,du-\mu_1^{(2)}\mu_0^{(2)}\}\).
The identical expansion with \(Q_1^{(2)}(1-u)\) gives
\[
\frac{N}{N-1}\left\{C_{\mathrm{ctr}}^{(2)}(F^{(2)})-(\theta_N^{(2)})^2\right\}
=\sum_b d_b^2S_b^{2,(2)}+2d_0d_1\underline S_{10}^{(2)}. \tag{16}
\]
Substitution of (15)--(16) into (13)--(14) proves
\[
U_N^{(2)}
=\sum_b\frac{d_b^2S_b^{2,(2)}}{m_b}
-\frac{\sum_b d_b^2S_b^{2,(2)}+2d_0d_1\overline S_{10}^{(2)}}{N},
\]
and
\[
L_N^{(2)}
=\sum_b\frac{d_b^2S_b^{2,(2)}}{m_b}
-\frac{\sum_b d_b^2S_b^{2,(2)}+2d_0d_1\underline S_{10}^{(2)}}{N}.
\]
These are the scaled-contrast versions of the sharp endpoints in Aronow, Green, and Lee (2014, Lemma~1 and equations~(7)--(11)); the argument above independently derives them from the stated complete-randomization law and the fixed empirical marginals.
